% ============================================================================
%  SECCIÓN 1.5: OBJETIVOS DE LA INVESTIGACIÓN
% ============================================================================

\section{Objetivos de la investigaci\'on}

Los objetivos de la investigaci\'on son el punto de referencia del estudio, ya
que orientan las acciones necesarias para alcanzar los resultados que se
pretenden lograr con el desarrollo del proyecto \citep{arias2012}. Deben
redactarse con verbos en infinitivo que implican acciones concretas que se
ejecutan durante la investigaci\'on.

\subsection{Objetivo general}

El objetivo general constituye el prop\'osito general que se persigue con el
proceso de investigaci\'on. Es un enunciado claro y preciso de la meta global,
expresado en t\'erminos cualitativos:

\begin{quote}
\textbf{Desarrollar} un sistema web y m\'ovil para la visualizaci\'on interactiva
de escenarios deportivos a nivel nacional que facilite a la ciudadan\'ia la
localizaci\'on y el acceso a informaci\'on confiable y actualizada sobre la
infraestructura deportiva del pa\'is.
\end{quote}

\subsection{Objetivos espec\'ificos}

Los objetivos espec\'ificos constituyen los prop\'ositos particulares mediante los
cuales se puede lograr el objetivo general e indican lo que se pretende
realizar en cada una de las etapas del proyecto. Para el presente proyecto se
han definido los siguientes objetivos espec\'ificos:

\begin{enumerate}
  \item \textbf{Diagnosticar} la situaci\'on actual de la difusi\'on y el acceso
    a la informaci\'on de los escenarios deportivos a nivel nacional, mediante
    t\'ecnicas de recolecci\'on de datos como entrevistas y encuestas
    \citep{hernandez2014}, a fin de fundamentar los requerimientos del
    sistema.

  \item \textbf{Identificar y documentar} los requerimientos funcionales y no
    funcionales del sistema web y m\'ovil, siguiendo las buenas pr\'acticas de la
    ingenier\'ia de requerimientos \citep{sommerville2011}.

  \item \textbf{Dise\~nar} la arquitectura del sistema, incluyendo el modelo de
    datos, la estructura de m\'odulos y la interfaz de usuario, mediante
    diagramas y wireframes de baja fidelidad.

  \item \textbf{Desarrollar} el m\'odulo de visualizaci\'on interactiva basado en
    mapas georreferenciados, que permita ubicar los escenarios deportivos y
    consultar su informaci\'on de manera intuitiva.

  \item \textbf{Implementar} la aplicaci\'on m\'ovil multiplataforma y el sistema
    web, integrando las funcionalidades definidas en el alcance del producto
    m\'inimo viable.

  \item \textbf{Evaluar} la usabilidad, el rendimiento y la satisfacci\'on de los
    usuarios con el sistema desarrollado, aplicando pruebas funcionales y
    cuestionarios, con el prop\'osito de validar el cumplimiento de los
    objetivos planteados.
\end{enumerate}
